\documentclass{article}
\usepackage[margin=25mm]{geometry}
\usepackage{amsmath,amssymb}
\usepackage{xurl}
\usepackage[hidelinks]{hyperref}
\setlength{\emergencystretch}{2em}
% Shared commands for notes in docs/paper-gaps/.

\DeclareUnicodeCharacter{2080}{\ifmmode{}_{0}\else\textsubscript{0}\fi}
\DeclareUnicodeCharacter{2081}{\ifmmode{}_{1}\else\textsubscript{1}\fi}
\DeclareUnicodeCharacter{2082}{\ifmmode{}_{2}\else\textsubscript{2}\fi}
\DeclareUnicodeCharacter{2099}{\ifmmode{}_{n}\else\textsubscript{n}\fi}

\newcommand{\Lean}{\textsc{Lean}}
\ExplSyntaxOn
\NewDocumentCommand{\leanid}{m}{
  \ifmmode
    \textsf{\detokenize{#1}}
  \else
    \group_begin:
    \urlstyle{sf}
    \tl_set:Nn \l_tmpa_tl {#1}
    \tl_replace_all:Nnn \l_tmpa_tl {\_} {_}
    \exp_args:NV \path \l_tmpa_tl
    \group_end:
  \fi
}
\ExplSyntaxOff
\providecommand{\tr}{\operatorname{tr}}
\newcommand{\tnleanIssueBase}{https://github.com/LionSR/TNLean/issues/}
\newcommand{\tnleanPrBase}{https://github.com/LionSR/TNLean/pull/}
\newcommand{\tnleanDocBase}{https://sirui-lu.com/TNLean/docs/}
\newcommand{\ghissue}[1]{\href{\tnleanIssueBase#1}{GitHub issue \##1}}
\newcommand{\ghpr}[1]{\href{\tnleanPrBase#1}{PR \##1}}
\newcommand{\doclink}[1]{\href{\tnleanDocBase#1.html}{\path{#1}}}

% Dirac notation and the identity operator, as in blueprint/src/macros/common.tex.
% \providecommand keeps note-local definitions authoritative.
\usepackage{dsfont}
\providecommand{\ket}[1]{|#1\rangle}
\providecommand{\bra}[1]{\langle #1|}
\providecommand{\braket}[2]{\langle #1 | #2 \rangle}
\providecommand{\ketbra}[2]{|#1\rangle\!\langle #2|}
\providecommand{\Id}{\mathds{1}}
\providecommand{\one}{\mathds{1}}
\providecommand{\C}{\mathbb{C}}
\providecommand{\MN}[1]{M_{#1}(\C)}
\providecommand{\MD}{M_D(\C)}

% Verdict marker: \gapnote{<kind>}{<status>}. Kind is the policy
% classification (clarification, local-correction, scope-restriction,
% unfaithful, false-source, open-gap); status is open, wip (elimination
% underway), resolved, or historical. Machine-read by the site index;
% prints a verdict line.
\newcommand{\gapnote}[2]{\par\noindent\textbf{Verdict:} #1 (#2).\par}

\title{AKLV's two exterior truncations: source conventions and large-window robustness}
\author{}
\date{2026-10-29}
\begin{document}
\maketitle
\gapnote{scope-restriction}{wip}

\section{The original bounded-term setting remains part of the argument}

The reference is AKLV, \href{https://arxiv.org/abs/1301.1162v1}{arXiv:1301.1162v1},
\path{1dfrustrated34.tex}. Lines 879--884 start with a nearest-neighbour sum
$H'=\sum_i h_i$ of positive contractions and a unique gapped ground vector.
Lines 890--899 introduce an exterior-block description. Lines 901--925 then
construct its exterior operators and scalar shifts from those original
bounded local terms. Theorem~6.1 at lines 949--963 refers to the resulting
$H$; it does not explicitly quantify over arbitrary independently supplied
exterior matrices.

This context is needed in the exponential part of the proof. Lines
1301--1312 expand the repeated commutators into the constituent local
interactions and use that their norms are at most one. A large norm for
an exterior \emph{sum} is compatible with this setting. An arbitrary
single-site matrix of the same large norm need not be such a sum.
The physical formalization retains the original positive-contraction
nearest-neighbour data. No counterexample to an arbitrary-exterior
extension is treated as a counterexample to this contextual source theorem.

\section{Two local corrections}

Line 897 assigns the left exterior to sites $1,\ldots,m$, whereas the
explicit definition at line 907 includes the bond $h_m$ on sites
$\{m,m+1\}$. The definition therefore includes the boundary site $m+1$.
The formalization follows the explicit original-bond sum. In zero-based
coordinates its left reference support is $[0,a+1)$, and adding the first
window bond enlarges it to $[0,a+2)$. This also gives the overlap needed
by the subsequent two-exterior truncation.

The inclusive middle range $4\le i\le s-3$ contains $s-6$ interactions,
whereas line 915 prints the individual subtraction $d/(s-7)$. For $s=8$
there are two interactions, so that prescription subtracts $2d$ from a
bulk with lowest energy $d$. If $d>0$, the resulting lowest energy is
$-d$ rather than zero. Subtracting the true minimum once from the aggregate
middle sum gives the intended normalization. Its equal redistribution
among the individual interactions uses $d/(s-6)$ when the range is nonempty.
No individual redistribution is needed for the proved aggregate identities.

\section{The actual normalization and two exterior clippings}

The three reference regions are
$[0,a+1)$, $[a+3,a+s-2)$, and $[a+s,L)$.
Their actual internal sums $A_j$ have pairwise disjoint supports. The
module \doclink{TNLean/MPS/AreaLaw/AKLV/Normalization} constructs genuine minima $e_j$
and a common unit vector attaining them. The remaining six boundary
contractions give an attained ground energy $E_0\in[0,6]$ for
$H=\sum_i h_i-(\sum_j e_j)\Id$.

For the three-region partition at $a+2,a+s-1$, its enlarged exterior sums
are exactly $C_0=A_0+h_a$ and $C_2=h_{a+s-1}+A_2$.
The two exterior clippings replace these by
$\min(C_0,(e_0+t)\Id)-e_0\Id$ and
$\min(C_2,(e_2+t)\Id)-e_2\Id$.
The middle sum and both connecting interactions are retained.
The minima of the reference sums are lower bounds for the enlarged sums;
they are not asserted to be their attained minima.

\section{A polynomial threshold and an aligned actual ground vector}

The module \doclink{TNLean/MPS/AreaLaw/AKLV/PhysicalRobustness} retains the actual unit
ground vector $\Omega$, energy $\mu$, and first excited energy $\mu+g$
of the original nearest-neighbour Hamiltonian, with $g>0$.
The ground-line gap and the original endpoint trial-vector argument derive
$g\le1$; this is not an independent hypothesis.
The genuine reference minima are chosen before $t$, with
$E_0+\sum_j e_j=\mu$ and $0\le E_0\le6$.
For
\begin{align*}
 t\ge\frac{1024}{g}\left(\frac{E_0}{g}+1\right),
\end{align*}
the actual two-exterior truncation has a phase-aligned unit ground vector
$\phi$ and actual spectral gap at least $9g/10$, with
\begin{align*}
 \|\phi-\Omega\|^2\le\frac{e^{-t/16}}{100}.
\end{align*}
This gives an explicit universal sufficient constant in the polynomial
threshold printed at source line 955, together with the exponential
conclusion at line 960. No perturbation, residual, truncated ground vector,
closeness estimate, or upper bound on the original gap is supplied.

\section{The stronger logarithmic threshold}

The same actual exponential residual comparison gives the stronger sufficient
condition
\begin{align*}
 t\ge T(g):=1024+16\log(1/g).
\end{align*}
Indeed, setting $\tau=e^{-t/16}$ in the spectral comparison requires only
\begin{align*}
 t&\ge g+11,&
 t&\ge2g+16\log(C_1/g),& C_1&=10^6\cdot64e^{3/4}.
\end{align*}
These follow from $0<g\le1$, $\log C_1\le61$, and $t\ge T(g)$.
The physical endpoint argument derives $g\le1$ from the original chain.
Thus the actual phase-aligned vector, gap and exponential error are available
above a logarithmic threshold independently of $E_0$.
The actual clipped Chebyshev filter, its spectral ceiling and sharp rank
bound are available at the same threshold. Since $T(g)\le1024/g$, the earlier
polynomial-threshold conclusions follow as corollaries without changing
their physical hypotheses. The later changing-target and entropy
construction remains separate.

\section{The remaining finite-window restriction}

The proved geometry assumes $s\ge6$, so the six boundary labels are
distinct; at $s=6$ the middle reference sum is empty and is normalized by
its genuine minimum zero. The source does not state a numerical minimum
for $s$ beside Theorem~6.1. Its application at line 1023 chooses
$s=O(\log^2 d/g)$ with a sufficiently large constant. The actual uniform
choice in \doclink{TNLean/MPS/AreaLaw/AKLV/FavorableFilter} proves $s\ge21$,
so this application satisfies the six-interaction restriction.
Windows with fewer than six interactions remain separate. The blueprint
states this restriction explicitly and does not identify the large-window
statement with an unrestricted all-window formulation of the printed theorem.
The actual favorable square-degree filter and initial product approximation
are constructed for every window selected by the universal scalar choice
that fits in the chain. The fixed-window changing-target sequence is now constructed by actual
normalized polynomial images with an exact product rank budget. Finite
small windows remain a separate robustness statement. The entropy
argument from the actual coherent sequence is completed below. The
accuracy-dependent enlarged-window optimization is completed below,
including exact boundary alternatives for chains too short for that window.
\section{The initial product overlap}

At source line 1023, the product-overlap inequality is printed with
$\le$; the cited Lemma~2.2 gives $\ge1/\sqrt{2D}$.
For the initial application to the clipped filter, its fixed vector is the
actual clipped ground vector $\phi^{t_0}$.
The construction in \doclink{TNLean/MPS/AreaLaw/AKLV/Initialization}
therefore chooses a product vector $\chi$ satisfying
\begin{align*}
 |\langle\phi^{t_0},\chi\rangle|\ge1/\sqrt{2D}.
\end{align*}
This makes the target explicit where the printed line writes $\Gamma$.
The robustness estimate alone does not transfer this lower bound to
$|\langle\Gamma,\chi\rangle|$: its error need not be small relative to
$D^{-1/2}$. No such transfer is used.
Instead, apply the actual clipped filter to $\chi$, normalize, and choose a
unit scalar phase. With squared perpendicular shrinking coefficient $\delta$,
the resulting unit vector $\psi_0$ satisfies
\begin{align*}
 \operatorname{SR}(\psi_0)&\le D,&
 \|\psi_0-\phi^{t_0}\|^2&\le4D\delta.
\end{align*}
The uniform parameter choice gives $D\delta\le10^{-5}$, so the same filter
already suffices for this initial approximation. The triangle inequality
then compares $\psi_0$ with $\Gamma$. The weaker bounds retained in the
statement are $\|\psi_0-\phi^{t_0}\|^2\le1/64$ and
$\|\psi_0-\Gamma\|^2\le1/16$.
This is the initial step of the changing-target construction, with the
available-window restriction stated separately; it is not the full
sequence or entropy theorem.

\section{The coherent changing-target sequence}

The construction in \doclink{TNLean/MPS/AreaLaw/AKLV/ChangingTargets}
uses one dependent recursion on the actual preceding unit vector. Its
fixed regional minima come from the physical initialization; the minimizing
energies returned by later refinements agree with them by uniqueness of
the minimum. Thus all thresholds truncate the same normalized physical
Hamiltonian at the same cut. The available-window and actual-first-excited-
energy restrictions remain explicit.

With $\beta_i=2^{-i}$ and
$t_i=1024+16\log(1/g)+16i\log2$, choose
$\ell_i=\lceil5\sqrt{(2t_i+s)/g}\rceil$ for $i\ge1$.
The proved sequence has $\|\psi_i-\Omega\|^2\le\beta_i/16$ and
rank at most the initial rank times the product of the sharp individual
filter bounds. Every edge retains its actual polynomial, phase and
normalized filtered image. This establishes a concrete sequence, rather
than independently selected approximants at different indices.

The printed summary at source line 1053 writes the rank itself as an
initial term plus a sum of logarithmic rank costs. The multiplicative
rank bound, or its logarithm, is the correct reading, consistent with
item 2 of the stated lemma. The current result retains the exact product
and the proved initial logarithmic budget. The finite entropy summation
from these actual ranks is completed below.
The accuracy-dependent enlarged-window argument at lines 1503--1517 is
also separate and is not replaced by this fixed-window recursion.

\section{Accuracy-dependent windows and the sharp rank comparison}

The construction in \doclink{TNLean/MPS/AreaLaw/AKLV/AccuracyWindow}
chooses the source's accuracy-dependent enlarged window at a prescribed
physical cut. A fitting enlarged window admits the actual initial window
at that same cut. Refinement applies the actual polynomial to that actual
initialized state. The two widths can have different regional minima;
no equality of their normalization Hamiltonians is assumed.
If the enlarged window does not fit, the ground vector itself is the
exact per-cut approximant, using the coefficient-rank bound
$q^{\min\{k,L-k\}}$. Thus endpoints and short chains need no supplied window.
The upper bound one for the actual gap is derived in the fitting case,
and is not a hypothesis of the all-cut statement.

For squared error $0<\rho\le1$, put $z=\log(128/\rho)$,
$m=\lceil\sqrt{\sqrt{z^3/g}}\rceil$, $s=s_0+2m$,
$t=T(g)+16z$, and $\ell=\lceil z\sqrt{(2t+s)/g}\rceil$.
The actual per-cut squared error is at most $\rho$. Its common rank budget
is the maximum of the exact filtered rank and exact boundary rank, with
logarithm
\begin{align*}
\max\bigl\{\log D_0+
32(\lfloor\ell/s\rfloor+\lfloor\sqrt\ell\rfloor+1)\log(q(\ell+1)),\,
(s-\lfloor s/2\rfloor)\log q\bigr\}.
\end{align*}
The construction in \doclink{TNLean/MPS/AreaLaw/AKLV/OptimizedRankBound}
now proves the uniform comparison
\begin{align*}
 \log B\le C_A\left[\frac{u^3}{g}+
      \left(\frac{z^3}{g}\right)^{1/4}(u+\log(z+1))\right].
\end{align*}
The universal coefficient depends only on the already universal $A$.
For $z\le s_0$, the actual degree is at most $66s_0^2$, and the extra
logarithmic rank cost is absorbed by the initial cost. For $s_0\le z$,
one has $m<z+1$, $2t+s\le4225z$, and the actual degree is at most
$66m^2$. The sharp polynomial-rank formula therefore has exponent at
most $2432m$. The exact boundary alternative has logarithmic cost at
most $s_0\log q+2m\log q$, which is absorbed by the same bound.
No exponential-in-degree replacement is used.

The physical theorem in
\doclink{TNLean/MPS/AreaLaw/AKLV/SharpPerCutApproximation}
derives $q\ge2$ and $g\le1$ from the two distinct spectral values and
the original unique-ground positive-contraction chain, independently
of window fit. Neither is a new hypothesis. The actual excited spectral
value $\mu+g$ remains explicit; the statement is not asserted for every
arbitrary projector-gap lower bound without choosing an actual gap.

Finally \doclink{TNLean/MPS/AreaLaw/AKLV/SharpMPSApproximation}
sets $\rho=\epsilon^2/[16(L+1)]$. The actual per-cut approximants give
Schmidt discarded mass at most $\rho$. Sequential truncation by the
leading marginal spectral projections has squared error at most
$\epsilon^2/16$ and cut ranks at most $B$. Its normalization is therefore
an actual unit matrix product state within norm error $\epsilon$, with
the same bond bound and
$z=\log(2048(L+1)/\epsilon^2)$. The returned vector is that normalized
ordered product. The proof uses no independent area-law or matrix
product approximation existence theorem. This completes this sharp
approximation route. The entropy summation from the separate coherent
changing-target sequence is also completed below.
\section{Entropy from the actual coherent sequence}

The module \doclink{TNLean/MPS/AreaLaw/AKLV/SequenceEntropyParameters}
proves, for the actual selected degrees and filter ranks,
\begin{align*}
 1\le\ell_j&\le240s(j+1),&
 \log F_j&\le200000(j+1)^2\log D_0.
\end{align*}
These are coarse estimates for entropy summation only; the separate
accuracy-dependent fourth-root rank comparison is unchanged.
The actual rank products satisfy $R_i\ge2^{i+1}$ and
$\log(R_{i+1}/R_i)=\log F_{i+1}$. The generic finite rank-tail entropy
inequality can therefore be applied with cutoff equal to the finite
marginal dimension. The existing finite dyadic quadratic moment gives
\begin{align*}
 \sum_{i<M}2^{-i}&\le2,&
 \sum_{i<M}2^{-i}(i+2)^2&\le22.
\end{align*}
The actual marginal tails $2^{-i}/16$ supplied by the coherent sequence
imply entropy at most $300000\log D_0$.

The module \doclink{TNLean/MPS/AreaLaw/AKLV/PhysicalEntropy} combines
this estimate with the actual initial-width bound. If the prescribed
initial window fails to fit, the exact smaller-side dimension estimate
$S_k\le\min\{k,L-k\}\log q\le s\log q$ suffices. The resulting theorem
holds at every cut, including the endpoints and chains shorter than the
initial window, with a universal constant chosen before the physical data:
\begin{align*}
 S_k(\Omega)\le\frac{C}{g}(1+\log q+\log(1/g))^3.
\end{align*}
It retains the actual first excited spectral value $\mu+g$. The local
dimension and upper-gap bounds are derived from the original physical
interaction independently of window fit. No supplied spectral tail,
approximant, filter, or independent area-law theorem is used. The displayed
logarithmic gap dependence is retained explicitly; no sharper dependence
with all inverse-gap logarithms removed is asserted.
\end{document}
